
\documentclass[11pt,a4paper]{article}
\usepackage[utf8]{inputenc}
\usepackage[T1]{fontenc}
\usepackage{lmodern}
\usepackage{microtype}
\usepackage{amsmath,amssymb,amsfonts,mathtools,bm}
\usepackage{physics}
\usepackage{geometry}
\usepackage{booktabs,longtable,array,multirow}
\usepackage{graphicx,float}
\usepackage{xcolor}
\usepackage{hyperref}
\usepackage{fancyhdr}
\usepackage{listings}
\usepackage{tcolorbox}
\usepackage{enumitem}
\usepackage{caption}
\geometry{margin=1.35cm}
\hypersetup{colorlinks=true, linkcolor=blue!55!black, urlcolor=blue!60!black, citecolor=blue!60!black}
\pagestyle{fancy}
\fancyhf{}
\lhead{NablaMath Professional LaTeX Engine}
\rhead{\leftmark}
\cfoot{\thepage}
\setlength{\parskip}{0.45em}
\setlength{\parindent}{0pt}
\lstdefinestyle{nabla}{
  basicstyle=\ttfamily\small,
  backgroundcolor=\color{black!3},
  frame=single,
  breaklines=true,
  keywordstyle=\color{blue!70!black},
  commentstyle=\color{green!40!black},
  stringstyle=\color{red!50!black},
  showstringspaces=false,
  columns=fullflexible
}
\newtcolorbox{derivationbox}{colback=blue!2!white,colframe=blue!45!black,title=Derivation}
\newtcolorbox{resultbox}{colback=green!2!white,colframe=green!45!black,title=Result}
\newtcolorbox{warningbox}{colback=orange!3!white,colframe=orange!60!black,title=Model scope}
\title{NablaMath Relativity Module - Complete Professional Derivation Report}
\author{NablaMath All-in-One Professional Engine}
\date{\today}
\begin{document}
\maketitle
\begin{abstract}
This document was generated by a single-file NablaMath professional module. The report intentionally uses real LaTeX for mathematical derivations, equations, tables, figures, appendices, and reproducibility metadata. It is designed as a template for producing long technical dossiers, including hundreds of pages when requested.
\end{abstract}
\tableofcontents
\newpage
\section{Objective and structure}
This module is designed to generate long-form relativity reports. It is not a short summary. It preserves assumptions, derives equations step by step, produces tables and figures, and exposes a configurable appendix generator for hundreds of pages.
\begin{warningbox}
The report is pedagogical and symbolic. It is not a substitute for a full numerical relativity code such as BSSN/ADM solvers, but it provides a rigorous derivation scaffold for special and general relativity.
\end{warningbox}

\part{Special Relativity}
\section{Postulates}
\begin{enumerate}[label=\textbf{P\arabic*.}]
\item The laws of physics are identical in all inertial frames.
\item The vacuum speed of light is invariant and equal to $c$ in all inertial frames.
\end{enumerate}
Let $S$ and $S'$ be inertial frames, with $S'$ moving at velocity $v$ along the $x$ axis relative to $S$.

\section{Linearity of transformations}
Homogeneity of space and time implies a linear coordinate map:
\begin{align}
x' &= A x + B t, \\
t' &= D x + E t.
\end{align}
The origin of $S'$ satisfies $x'=0$ and moves as $x=vt$. Hence
\begin{align}
0=A(vt)+Bt \Rightarrow B=-Av,
\end{align}
giving
\begin{align}
x'=A(x-vt).
\end{align}
By symmetry the inverse must have the same form with $v\to -v$:
\begin{align}
x=A(x'+vt').
\end{align}

\section{Light invariance and Lorentz factor}
Light moving along $+x$ obeys $x=ct$ and $x'=ct'$. Light moving along $-x$ obeys $x=-ct$ and $x'=-ct'$. The interval condition is
\begin{align}
c^2t'^2-x'^2=c^2t^2-x^2.
\end{align}
Substitute the ansatz $x'=A(x-vt)$ and solve for $t'$ by requiring interval preservation:
\begin{derivationbox}
\begin{align}
x'&=\gamma(x-vt),\\
ct'&=\gamma\left(ct-\frac{v}{c}x\right),\\
t'&=\gamma\left(t-\frac{vx}{c^2}\right),\\
\gamma&=\frac{1}{\sqrt{1-v^2/c^2}}.
\end{align}
\end{derivationbox}
Therefore the Lorentz transformation is
\begin{resultbox}
\begin{align}
x'&=\gamma(x-vt),\\
y'&=y,\\
z'&=z,\\
t'&=\gamma\left(t-\frac{vx}{c^2}\right).
\end{align}
\end{resultbox}

\section{Time dilation}
A clock at rest in $S'$ has $\Delta x'=0$. From inverse Lorentz transformation,
\begin{align}
\Delta t=\gamma\left(\Delta t'+\frac{v\Delta x'}{c^2}\right)=\gamma\Delta t'.
\end{align}
If $\Delta\tau=\Delta t'$ is proper time, then
\begin{resultbox}
\[
\Delta t=\gamma\Delta\tau.
\]
\end{resultbox}

\section{Length contraction}
A rod at rest in $S'$ has proper length $L_0=\Delta x'$. Its length in $S$ is measured simultaneously, so $\Delta t=0$:
\begin{align}
\Delta x' &= \gamma(\Delta x-v\Delta t)=\gamma\Delta x.
\end{align}
Thus
\begin{resultbox}
\[
L=\Delta x=\frac{L_0}{\gamma}.
\]
\end{resultbox}

\section{Relativistic velocity addition}
Let $u=dx/dt$ and $u'=dx'/dt'$. Then
\begin{align}
u' &= \frac{dx'}{dt'}
=\frac{\gamma(dx-vdt)}{\gamma(dt-vdx/c^2)}
=\frac{u-v}{1-uv/c^2}.
\end{align}
The inverse is
\[
u=\frac{u'+v}{1+u'v/c^2}.
\]

\section{Four-vectors and invariant interval}
Define the spacetime coordinate
\[
x^\mu=(ct,x,y,z).
\]
With metric signature $(-,+,+,+)$,
\begin{align}
ds^2 &= \eta_{\mu\nu}dx^\mu dx^\nu \\
     &= -c^2dt^2+dx^2+dy^2+dz^2.
\end{align}
Timelike proper time satisfies
\begin{align}
d\tau^2 = dt^2-\frac{1}{c^2}(dx^2+dy^2+dz^2).
\end{align}

\section{Energy and momentum}
Four-velocity is
\begin{align}
U^\mu=\dv{x^\mu}{\tau}=\gamma(c,\mathbf{v}).
\end{align}
Four-momentum is
\begin{align}
p^\mu=mU^\mu=(\gamma mc,\gamma m\mathbf{v}).
\end{align}
Define
\begin{align}
E=\gamma mc^2,\qquad \mathbf{p}=\gamma m\mathbf{v}.
\end{align}
The invariant norm becomes
\begin{align}
p_\mu p^\mu&=-m^2c^2,\\
-\frac{E^2}{c^2}+p^2&=-m^2c^2,\\
E^2&=p^2c^2+m^2c^4.
\end{align}
At rest $p=0$, so
\begin{resultbox}
\[
E_0=mc^2.
\]
\end{resultbox}

\begin{figure}[H]
\centering
\includegraphics[width=0.92\textwidth]{lorentz_factor.png}
\caption{Lorentz factor growth as velocity approaches c.}
\end{figure}

\begin{figure}[H]
\centering
\includegraphics[width=0.65\textwidth]{light_cone.png}
\caption{Causal structure of flat spacetime in 1+1 dimensions.}
\end{figure}

\part{General Relativity}
\section{Equivalence principle}
Locally, a uniform gravitational field is indistinguishable from acceleration. This motivates replacing gravitational force by geometry: free particles follow geodesics in curved spacetime.

\section{Metric tensor}
The invariant interval generalizes from
\[
ds^2=\eta_{\mu\nu}dx^\mu dx^\nu
\]
to
\begin{resultbox}
\[
ds^2=g_{\mu\nu}(x)dx^\mu dx^\nu.
\]
\end{resultbox}
The metric contains local clock rates, rulers, causal cones, and gravitational potentials.

\section{Christoffel symbols}
The Levi-Civita connection is torsion-free and metric-compatible:
\begin{align}
\nabla_\lambda g_{\mu\nu}=0,\qquad \Gamma^\rho_{\mu\nu}=\Gamma^\rho_{\nu\mu}.
\end{align}
Solving these conditions gives
\begin{derivationbox}
\[
\Gamma^\rho_{\mu\nu}=\frac{1}{2}g^{\rho\sigma}
\left(\partial_\mu g_{\nu\sigma}+\partial_\nu g_{\mu\sigma}-\partial_\sigma g_{\mu\nu}\right).
\]
\end{derivationbox}

\section{Geodesic equation from extremal proper time}
The action for a free massive particle is
\begin{align}
S=-mc\int ds=-mc\int \sqrt{-g_{\mu\nu}\dot x^\mu\dot x^\nu}\,d\lambda.
\end{align}
The Euler-Lagrange equations lead to
\begin{resultbox}
\[
\dv[2]{x^\rho}{\tau}+\Gamma^\rho_{\mu\nu}\dv{x^\mu}{\tau}\dv{x^\nu}{\tau}=0.
\]
\end{resultbox}

\section{Curvature tensor}
Curvature measures the non-commutation of covariant derivatives:
\begin{align}
[\nabla_\mu,\nabla_\nu]V^\rho=R^\rho{}_{\sigma\mu\nu}V^\sigma.
\end{align}
In terms of Christoffel symbols,
\begin{align}
R^\rho{}_{\sigma\mu\nu}
&=\partial_\mu\Gamma^\rho_{\nu\sigma}
-\partial_\nu\Gamma^\rho_{\mu\sigma}
+\Gamma^\rho_{\mu\lambda}\Gamma^\lambda_{\nu\sigma}
-\Gamma^\rho_{\nu\lambda}\Gamma^\lambda_{\mu\sigma}.
\end{align}
Contracting indices gives the Ricci tensor and scalar:
\begin{align}
R_{\mu\nu}=R^\rho{}_{\mu\rho\nu},\qquad R=g^{\mu\nu}R_{\mu\nu}.
\end{align}

\section{Einstein-Hilbert action}
The gravitational action is
\begin{align}
S_g=\frac{c^3}{16\pi G}\int R\sqrt{-g}\,d^4x.
\end{align}
Add matter action $S_m$ and vary with respect to $g^{\mu\nu}$:
\begin{align}
\delta(S_g+S_m)=0.
\end{align}
The variation gives
\begin{resultbox}
\[
R_{\mu\nu}-\frac{1}{2}Rg_{\mu\nu}+\Lambda g_{\mu\nu}=\frac{8\pi G}{c^4}T_{\mu\nu}.
\]
\end{resultbox}

\section{Newtonian limit}
In weak, static fields,
\begin{align}
g_{00}\approx -\left(1+\frac{2\Phi}{c^2}\right).
\end{align}
The geodesic equation reduces to
\begin{align}
\dv[2]{x^i}{t}\approx -\partial_i\Phi.
\end{align}
The $00$ component of Einstein's equation reduces to Poisson's equation:
\begin{align}
\nabla^2\Phi=4\pi G\rho.
\end{align}

\section{Schwarzschild metric}
For a static, spherically symmetric vacuum exterior,
\begin{resultbox}
\[
ds^2=-\left(1-\frac{2GM}{rc^2}\right)c^2dt^2+
\left(1-\frac{2GM}{rc^2}\right)^{-1}dr^2+r^2d\Omega^2.
\]
\end{resultbox}
Define Schwarzschild radius
\[
r_s=\frac{2GM}{c^2}.
\]
Gravitational redshift for a photon emitted at radius $r$ and received at infinity is
\begin{align}
1+z=\left(1-\frac{r_s}{r}\right)^{-1/2}.
\end{align}

\section{Perihelion precession}
The weak-field precession per orbit is
\begin{resultbox}
\[
\Delta\varpi=\frac{6\pi GM}{a(1-e^2)c^2}.
\]
\end{resultbox}

\begin{figure}[H]
\centering
\includegraphics[width=0.92\textwidth]{schwarzschild_redshift.png}
\caption{Schwarzschild gravitational redshift.}
\end{figure}

\begin{figure}[H]
\centering
\includegraphics[width=0.65\textwidth]{precessing_orbit.png}
\caption{Schematic relativistic periapsis precession.}
\end{figure}

\part{Extended calculation ledger and appendices}

\section{Relativity ledger case 1: $\beta=0.0500$}
This ledger entry expands the same derivation numerically, showing how the symbolic equations become computable quantities.
\begin{align}
\beta &= \frac{v}{c}=0.050000,\\
v &= \beta c = 1.4990e+07\,\mathrm{m/s},\\
\gamma &= \frac{1}{\sqrt{1-\beta^2}} = 1.00125235,\\
\Delta t &= \gamma\Delta\tau = 1.00125235(1.0\,\mathrm{s}) = 1.00125235\,\mathrm{s},\\
L &= \frac{L_0}{\gamma} = \frac{10.0\,\mathrm{m}}{1.00125235} = 9.98749218\,\mathrm{m},\\
\frac{E}{mc^2}&=\gamma=1.00125235.
\end{align}

For a Schwarzschild redshift calculation at $r=2.6200r_s$:
\begin{align}
r_s &= \frac{2GM_\odot}{c^2}=2953.339382\,\mathrm{m},\\
1+z &= \left(1-\frac{r_s}{r}\right)^{-1/2}
=\left(1-\frac{1}{2.6200}\right)^{-1/2},\\
z&=0.27172479.
\end{align}

\section{Relativity ledger case 2: $\beta=0.1175$}
This ledger entry expands the same derivation numerically, showing how the symbolic equations become computable quantities.
\begin{align}
\beta &= \frac{v}{c}=0.117500,\\
v &= \beta c = 3.5226e+07\,\mathrm{m/s},\\
\gamma &= \frac{1}{\sqrt{1-\beta^2}} = 1.00697544,\\
\Delta t &= \gamma\Delta\tau = 1.00697544(1.0\,\mathrm{s}) = 1.00697544\,\mathrm{s},\\
L &= \frac{L_0}{\gamma} = \frac{10.0\,\mathrm{m}}{1.00697544} = 9.93072883\,\mathrm{m},\\
\frac{E}{mc^2}&=\gamma=1.00697544.
\end{align}

For a Schwarzschild redshift calculation at $r=2.7400r_s$:
\begin{align}
r_s &= \frac{2GM_\odot}{c^2}=2953.339382\,\mathrm{m},\\
1+z &= \left(1-\frac{r_s}{r}\right)^{-1/2}
=\left(1-\frac{1}{2.7400}\right)^{-1/2},\\
z&=0.25487555.
\end{align}

\newpage

\section{Relativity ledger case 3: $\beta=0.1850$}
This ledger entry expands the same derivation numerically, showing how the symbolic equations become computable quantities.
\begin{align}
\beta &= \frac{v}{c}=0.185000,\\
v &= \beta c = 5.5462e+07\,\mathrm{m/s},\\
\gamma &= \frac{1}{\sqrt{1-\beta^2}} = 1.01756467,\\
\Delta t &= \gamma\Delta\tau = 1.01756467(1.0\,\mathrm{s}) = 1.01756467\,\mathrm{s},\\
L &= \frac{L_0}{\gamma} = \frac{10.0\,\mathrm{m}}{1.01756467} = 9.82738521\,\mathrm{m},\\
\frac{E}{mc^2}&=\gamma=1.01756467.
\end{align}

For a Schwarzschild redshift calculation at $r=2.8600r_s$:
\begin{align}
r_s &= \frac{2GM_\odot}{c^2}=2953.339382\,\mathrm{m},\\
1+z &= \left(1-\frac{r_s}{r}\right)^{-1/2}
=\left(1-\frac{1}{2.8600}\right)^{-1/2},\\
z&=0.24001387.
\end{align}

\section{Relativity ledger case 4: $\beta=0.2525$}
This ledger entry expands the same derivation numerically, showing how the symbolic equations become computable quantities.
\begin{align}
\beta &= \frac{v}{c}=0.252500,\\
v &= \beta c = 7.5698e+07\,\mathrm{m/s},\\
\gamma &= \frac{1}{\sqrt{1-\beta^2}} = 1.03348823,\\
\Delta t &= \gamma\Delta\tau = 1.03348823(1.0\,\mathrm{s}) = 1.03348823\,\mathrm{s},\\
L &= \frac{L_0}{\gamma} = \frac{10.0\,\mathrm{m}}{1.03348823} = 9.67596894\,\mathrm{m},\\
\frac{E}{mc^2}&=\gamma=1.03348823.
\end{align}

For a Schwarzschild redshift calculation at $r=2.9800r_s$:
\begin{align}
r_s &= \frac{2GM_\odot}{c^2}=2953.339382\,\mathrm{m},\\
1+z &= \left(1-\frac{r_s}{r}\right)^{-1/2}
=\left(1-\frac{1}{2.9800}\right)^{-1/2},\\
z&=0.22680500.
\end{align}

\newpage

\section{Relativity ledger case 5: $\beta=0.3200$}
This ledger entry expands the same derivation numerically, showing how the symbolic equations become computable quantities.
\begin{align}
\beta &= \frac{v}{c}=0.320000,\\
v &= \beta c = 9.5934e+07\,\mathrm{m/s},\\
\gamma &= \frac{1}{\sqrt{1-\beta^2}} = 1.05550083,\\
\Delta t &= \gamma\Delta\tau = 1.05550083(1.0\,\mathrm{s}) = 1.05550083\,\mathrm{s},\\
L &= \frac{L_0}{\gamma} = \frac{10.0\,\mathrm{m}}{1.05550083} = 9.47417543\,\mathrm{m},\\
\frac{E}{mc^2}&=\gamma=1.05550083.
\end{align}

For a Schwarzschild redshift calculation at $r=3.1000r_s$:
\begin{align}
r_s &= \frac{2GM_\odot}{c^2}=2953.339382\,\mathrm{m},\\
1+z &= \left(1-\frac{r_s}{r}\right)^{-1/2}
=\left(1-\frac{1}{3.1000}\right)^{-1/2},\\
z&=0.21498579.
\end{align}

\section{Relativity ledger case 6: $\beta=0.3875$}
This ledger entry expands the same derivation numerically, showing how the symbolic equations become computable quantities.
\begin{align}
\beta &= \frac{v}{c}=0.387500,\\
v &= \beta c = 1.1617e+08\,\mathrm{m/s},\\
\gamma &= \frac{1}{\sqrt{1-\beta^2}} = 1.08475200,\\
\Delta t &= \gamma\Delta\tau = 1.08475200(1.0\,\mathrm{s}) = 1.08475200\,\mathrm{s},\\
L &= \frac{L_0}{\gamma} = \frac{10.0\,\mathrm{m}}{1.08475200} = 9.21869703\,\mathrm{m},\\
\frac{E}{mc^2}&=\gamma=1.08475200.
\end{align}

For a Schwarzschild redshift calculation at $r=3.2200r_s$:
\begin{align}
r_s &= \frac{2GM_\odot}{c^2}=2953.339382\,\mathrm{m},\\
1+z &= \left(1-\frac{r_s}{r}\right)^{-1/2}
=\left(1-\frac{1}{3.2200}\right)^{-1/2},\\
z&=0.20434648.
\end{align}

\newpage

\section{Relativity ledger case 7: $\beta=0.4550$}
This ledger entry expands the same derivation numerically, showing how the symbolic equations become computable quantities.
\begin{align}
\beta &= \frac{v}{c}=0.455000,\\
v &= \beta c = 1.3641e+08\,\mathrm{m/s},\\
\gamma &= \frac{1}{\sqrt{1-\beta^2}} = 1.12297542,\\
\Delta t &= \gamma\Delta\tau = 1.12297542(1.0\,\mathrm{s}) = 1.12297542\,\mathrm{s},\\
L &= \frac{L_0}{\gamma} = \frac{10.0\,\mathrm{m}}{1.12297542} = 8.90491437\,\mathrm{m},\\
\frac{E}{mc^2}&=\gamma=1.12297542.
\end{align}

For a Schwarzschild redshift calculation at $r=3.3400r_s$:
\begin{align}
r_s &= \frac{2GM_\odot}{c^2}=2953.339382\,\mathrm{m},\\
1+z &= \left(1-\frac{r_s}{r}\right)^{-1/2}
=\left(1-\frac{1}{3.3400}\right)^{-1/2},\\
z&=0.19471772.
\end{align}

\section{Relativity ledger case 8: $\beta=0.5225$}
This ledger entry expands the same derivation numerically, showing how the symbolic equations become computable quantities.
\begin{align}
\beta &= \frac{v}{c}=0.522500,\\
v &= \beta c = 1.5664e+08\,\mathrm{m/s},\\
\gamma &= \frac{1}{\sqrt{1-\beta^2}} = 1.17282891,\\
\Delta t &= \gamma\Delta\tau = 1.17282891(1.0\,\mathrm{s}) = 1.17282891\,\mathrm{s},\\
L &= \frac{L_0}{\gamma} = \frac{10.0\,\mathrm{m}}{1.17282891} = 8.52639285\,\mathrm{m},\\
\frac{E}{mc^2}&=\gamma=1.17282891.
\end{align}

For a Schwarzschild redshift calculation at $r=3.4600r_s$:
\begin{align}
r_s &= \frac{2GM_\odot}{c^2}=2953.339382\,\mathrm{m},\\
1+z &= \left(1-\frac{r_s}{r}\right)^{-1/2}
=\left(1-\frac{1}{3.4600}\right)^{-1/2},\\
z&=0.18596124.
\end{align}

\newpage

\section{Relativity ledger case 9: $\beta=0.5900$}
This ledger entry expands the same derivation numerically, showing how the symbolic equations become computable quantities.
\begin{align}
\beta &= \frac{v}{c}=0.590000,\\
v &= \beta c = 1.7688e+08\,\mathrm{m/s},\\
\gamma &= \frac{1}{\sqrt{1-\beta^2}} = 1.23853850,\\
\Delta t &= \gamma\Delta\tau = 1.23853850(1.0\,\mathrm{s}) = 1.23853850\,\mathrm{s},\\
L &= \frac{L_0}{\gamma} = \frac{10.0\,\mathrm{m}}{1.23853850} = 8.07403245\,\mathrm{m},\\
\frac{E}{mc^2}&=\gamma=1.23853850.
\end{align}

For a Schwarzschild redshift calculation at $r=3.5800r_s$:
\begin{align}
r_s &= \frac{2GM_\odot}{c^2}=2953.339382\,\mathrm{m},\\
1+z &= \left(1-\frac{r_s}{r}\right)^{-1/2}
=\left(1-\frac{1}{3.5800}\right)^{-1/2},\\
z&=0.17796303.
\end{align}

\section{Relativity ledger case 10: $\beta=0.6575$}
This ledger entry expands the same derivation numerically, showing how the symbolic equations become computable quantities.
\begin{align}
\beta &= \frac{v}{c}=0.657500,\\
v &= \beta c = 1.9711e+08\,\mathrm{m/s},\\
\gamma &= \frac{1}{\sqrt{1-\beta^2}} = 1.32722008,\\
\Delta t &= \gamma\Delta\tau = 1.32722008(1.0\,\mathrm{s}) = 1.32722008\,\mathrm{s},\\
L &= \frac{L_0}{\gamma} = \frac{10.0\,\mathrm{m}}{1.32722008} = 7.53454544\,\mathrm{m},\\
\frac{E}{mc^2}&=\gamma=1.32722008.
\end{align}

For a Schwarzschild redshift calculation at $r=3.7000r_s$:
\begin{align}
r_s &= \frac{2GM_\odot}{c^2}=2953.339382\,\mathrm{m},\\
1+z &= \left(1-\frac{r_s}{r}\right)^{-1/2}
=\left(1-\frac{1}{3.7000}\right)^{-1/2},\\
z&=0.17062819.
\end{align}

\newpage

\section{Relativity ledger case 11: $\beta=0.7250$}
This ledger entry expands the same derivation numerically, showing how the symbolic equations become computable quantities.
\begin{align}
\beta &= \frac{v}{c}=0.725000,\\
v &= \beta c = 2.1735e+08\,\mathrm{m/s},\\
\gamma &= \frac{1}{\sqrt{1-\beta^2}} = 1.45190802,\\
\Delta t &= \gamma\Delta\tau = 1.45190802(1.0\,\mathrm{s}) = 1.45190802\,\mathrm{s},\\
L &= \frac{L_0}{\gamma} = \frac{10.0\,\mathrm{m}}{1.45190802} = 6.88748866\,\mathrm{m},\\
\frac{E}{mc^2}&=\gamma=1.45190802.
\end{align}

For a Schwarzschild redshift calculation at $r=3.8200r_s$:
\begin{align}
r_s &= \frac{2GM_\odot}{c^2}=2953.339382\,\mathrm{m},\\
1+z &= \left(1-\frac{r_s}{r}\right)^{-1/2}
=\left(1-\frac{1}{3.8200}\right)^{-1/2},\\
z&=0.16387711.
\end{align}

\section{Relativity ledger case 12: $\beta=0.7925$}
This ledger entry expands the same derivation numerically, showing how the symbolic equations become computable quantities.
\begin{align}
\beta &= \frac{v}{c}=0.792500,\\
v &= \beta c = 2.3759e+08\,\mathrm{m/s},\\
\gamma &= \frac{1}{\sqrt{1-\beta^2}} = 1.63968856,\\
\Delta t &= \gamma\Delta\tau = 1.63968856(1.0\,\mathrm{s}) = 1.63968856\,\mathrm{s},\\
L &= \frac{L_0}{\gamma} = \frac{10.0\,\mathrm{m}}{1.63968856} = 6.09871913\,\mathrm{m},\\
\frac{E}{mc^2}&=\gamma=1.63968856.
\end{align}

For a Schwarzschild redshift calculation at $r=3.9400r_s$:
\begin{align}
r_s &= \frac{2GM_\odot}{c^2}=2953.339382\,\mathrm{m},\\
1+z &= \left(1-\frac{r_s}{r}\right)^{-1/2}
=\left(1-\frac{1}{3.9400}\right)^{-1/2},\\
z&=0.15764246.
\end{align}

\newpage

\section{Relativity ledger case 13: $\beta=0.8600$}
This ledger entry expands the same derivation numerically, showing how the symbolic equations become computable quantities.
\begin{align}
\beta &= \frac{v}{c}=0.860000,\\
v &= \beta c = 2.5782e+08\,\mathrm{m/s},\\
\gamma &= \frac{1}{\sqrt{1-\beta^2}} = 1.95965450,\\
\Delta t &= \gamma\Delta\tau = 1.95965450(1.0\,\mathrm{s}) = 1.95965450\,\mathrm{s},\\
L &= \frac{L_0}{\gamma} = \frac{10.0\,\mathrm{m}}{1.95965450} = 5.10294033\,\mathrm{m},\\
\frac{E}{mc^2}&=\gamma=1.95965450.
\end{align}

For a Schwarzschild redshift calculation at $r=4.0600r_s$:
\begin{align}
r_s &= \frac{2GM_\odot}{c^2}=2953.339382\,\mathrm{m},\\
1+z &= \left(1-\frac{r_s}{r}\right)^{-1/2}
=\left(1-\frac{1}{4.0600}\right)^{-1/2},\\
z&=0.15186691.
\end{align}

\section{Relativity ledger case 14: $\beta=0.9275$}
This ledger entry expands the same derivation numerically, showing how the symbolic equations become computable quantities.
\begin{align}
\beta &= \frac{v}{c}=0.927500,\\
v &= \beta c = 2.7806e+08\,\mathrm{m/s},\\
\gamma &= \frac{1}{\sqrt{1-\beta^2}} = 2.67506170,\\
\Delta t &= \gamma\Delta\tau = 2.67506170(1.0\,\mathrm{s}) = 2.67506170\,\mathrm{s},\\
L &= \frac{L_0}{\gamma} = \frac{10.0\,\mathrm{m}}{2.67506170} = 3.73823153\,\mathrm{m},\\
\frac{E}{mc^2}&=\gamma=2.67506170.
\end{align}

For a Schwarzschild redshift calculation at $r=4.1800r_s$:
\begin{align}
r_s &= \frac{2GM_\odot}{c^2}=2953.339382\,\mathrm{m},\\
1+z &= \left(1-\frac{r_s}{r}\right)^{-1/2}
=\left(1-\frac{1}{4.1800}\right)^{-1/2},\\
z&=0.14650138.
\end{align}

\newpage

\section{Relativity ledger case 15: $\beta=0.9950$}
This ledger entry expands the same derivation numerically, showing how the symbolic equations become computable quantities.
\begin{align}
\beta &= \frac{v}{c}=0.995000,\\
v &= \beta c = 2.9829e+08\,\mathrm{m/s},\\
\gamma &= \frac{1}{\sqrt{1-\beta^2}} = 10.01252349,\\
\Delta t &= \gamma\Delta\tau = 10.01252349(1.0\,\mathrm{s}) = 10.01252349\,\mathrm{s},\\
L &= \frac{L_0}{\gamma} = \frac{10.0\,\mathrm{m}}{10.01252349} = 0.99874922\,\mathrm{m},\\
\frac{E}{mc^2}&=\gamma=10.01252349.
\end{align}

For a Schwarzschild redshift calculation at $r=4.3000r_s$:
\begin{align}
r_s &= \frac{2GM_\odot}{c^2}=2953.339382\,\mathrm{m},\\
1+z &= \left(1-\frac{r_s}{r}\right)^{-1/2}
=\left(1-\frac{1}{4.3000}\right)^{-1/2},\\
z&=0.14150353.
\end{align}

\end{document}